\begin{teilaufgaben}
\item
The function $u(x,t)$ is defined on the strip
$\Omega = [0,3] \times [0,\infty)$ and
satisfies in $\Omega$ the equation
\[
u_{t}(x,t) = u_{xx}(x,t)
\]
and  $u(0,t) = u(3,t) = 1$ as well as
\[
u(x,0) = 1 + 3 \cdot x - x^2
\]
on the boundary of $\Omega$.

Determine approximate values for $u(1,2)$ and $u(2,2)$.
Use the explicit method of Richardson with $\Delta x = 1$ and $\Delta t = 1$.

\item
The function $u(x,t)$ is defined on the upper half plane
$\Omega = (-\infty, \infty) \times [0,\infty)$ and
satisfies in $\Omega$ the equation
\[
u_{x}(x,t) + u_{t}(x,t) = 0
\]
and $u(x,0) = |\cos(\pi/2 \cdot x)|$ on the boundary of $\Omega$. 

Determine approximate values for $u(-1,1)$, $u(0,1)$ and $u(1,1)$.
Use the scheme of Lax-Wendroff with $\Delta x = 1$ and $\Delta t = 1/2.$
\end{teilaufgaben}

\begin{loesung}
\begin{teilaufgaben}
\item
We apply the explicit method of Richardson with $r = 1$. 

We obtain with $k = 1$ and $2$ and $j = 0$ and $1$ the following iteration
formula
\begin{equation}
\tilde u(k,j+1)
=
\tilde u(k-1,j) - \tilde u(k,j) + \tilde u(k+1,j).
\tag{\bf{1P}}
\end{equation}
Therefore we have for the first time-step:
\[
\tilde u(1,1)
=
1 - \tilde u(1,0) + \tilde u(2,0).
\]
and
\[
\tilde u(2,1)
=
\tilde u(1,0) - \tilde u(2,0) + 1.
\]
With
\[
\tilde u(1,0) = \tilde u(2,0)
=
3.
\]
we obtain
\begin{equation}
\tilde u(1,1) = \tilde u(2,1) = 1.
\tag{\bf{1P}}
\end{equation}
and in a similar way
\begin{equation}
\tilde u(1,2) = \tilde u(2,2) = 1.
\tag{\bf{1P}}
\end{equation}
\item
We apply the Lax-Wendroff scheme with $r = 1/2$. 

We obtain the following iteration formula for $j = 0$ and $1/2$
\begin{equation}
\tilde u(k, j+1/2)
=
3/8 \cdot \tilde u(k-1,j) + 6/8 \cdot \tilde u(k,j) - 1/8
\cdot
\tilde u(k+1,j).
\tag{\bf{1P}}
\end{equation}
We start with
\[
\tilde u(*,0) = (0, 1, 0, 1, 0, 1, 0),
\]
where the first index ranges from $-3$ to $3$ and obtain
\begin{equation}
\tilde u(*,1/2) = (3/4, 1/4, 3/4, 1/4, 3/4),
\tag{\bf{1P}}
\end{equation}
where the first index ranges from $-2$ to $2$.
The second iteration leads to
\begin{equation}
\big(\tilde u(-1,1), \tilde u(0,1), \tilde u(1,1)\big)
=
(3/8, 5/8, 3/8).
\tag{\bf{1P}}
\end{equation}
\end{teilaufgaben}
\end{loesung}
